% !PRINTED none

\begin{center}
\gesperrt{Seventeenth Chapter.}

\gesperrt{Theory of Algebraic Fields.}

\rule{4cm}{0.4pt}

\S.\,144.

\gesperrt{Basis of an algebraic number field.}

\gesperrt{Discriminants.}
\end{center}

Let
\weq{1}{f(\Theta) = \Theta^n + a_1\Theta^{n-1} + \cdots + a_n = 0}
be the irreducible equation of degree $n$ with rational coefficients $a_1, a_2, \ldots, a_n$, which defines for us an algebraic field $\Omega = R(\Theta)$. The field $\Omega$ is then the totality of all numbers of the form
\weq{2}{\omega = h_1 + h_2\Theta + h_3\Theta^2 + \cdots + h_n\Theta^{n-1},}
in which $h_1, h_2, \ldots, h_n$ are \gesperrt{rational numbers}. But if we put rational functionals for $h_1, h_2, \ldots, h_n$, we obtain from (2) all \gesperrt{functionals} of the field $\Omega$.

If we consider an arbitrary system of $n$ numbers
\weq{3}{\omega_r = h_{1,r} + h_{2,r}\Theta + \ldots + h_{n,r}\Theta^{n-1}, \quad r = 1,2,\ldots,n}
under the assumption that the determinant
\weq{4}{H = \Sigma \pm h_{1,1}h_{2,2}\ldots h_{n,n}}
is different from zero, then, on account of the irreducibility of $f$, an equation of the form
\weq{5}{k_1\omega_1 + k_2\omega_2 + \cdots + k_n\omega_n = 0,}
in which $k_1, k_2, \ldots, k_n$ are rational numbers, can subsist only when these coefficients are all zero, and this still holds even when \gesperrt{rational functionals} are admitted for the coefficients in equation (5).
